
\subsection{The elliptic covering functors with their full framings}

\begin{result}{B.29}{Proved}{Cubic and quartic sphere models, half turns and the elliptic cut change}
The two elliptic pieces in B.27 admit explicit finite presentations
obtained from three- and four-boundary sphere moves, finite
relabeling quotients, and the central framing maps below.  The
covering functors are induced by the actual cubic and quartic
families; they have inverse functors and preserve every specified
boundary half turn.  Their mixed disk presentations have counts
\[
\begin{array}{c|rrrr}
 &\text{objects}&F\text{ arrows}&B\text{ arrows}&\text{relators}\\\hline
 B_3&2&1&2&1\\
 B_4&5&5&6&7.
\end{array}                                                     \tag{B.88}
\]
These presentations, all comparison words, and the identities used
in the boundary flags are computed in
\texttt{b\_sphere\_covers.py}.  In particular, the elliptic change
of cut is the explicit word $s=aba$; its framing relations are
\[
 s^2=(ab)^3,\qquad s^4=z\quad\text{in }H_{1,1}.                 \tag{B.89}
\]
The central lattice used in this cubic comparison sends the
elliptic full turn to the \emph{square} of the disk full twist.
The comparison therefore includes a specified change of framing
lattice, rather than identifying those two turns.
\end{result}
\begin{proof}
First retain the outer boundary in B.28: remove its single
outer-capping relator.  The resulting groupoid is the fundamental
groupoid of the ordered configuration space of $m$ points in a
disk, for $m=3,4$, with the stated parenthesized markings.  The
forward and inverse formulas (B.84)--(B.85) remain valid; their
proof simply omits the last relator.  Permuting the point labels
acts freely on its objects $(p,T)$ and commutes with both word
maps.  Its ordinary orbit groupoid has objects the shapes $T$.
Every orbit arrow has a unique lift after its starting ordering
is fixed, so the orbit functor is the actual covering functor
from ordered to unordered configurations.

There are two shapes at $m=3$: one $F$ edge and one cherry
half turn at each shape.  Contracting the $F$ edge leaves the
two adjacent braid generators and their one braid relation.
At $m=4$ there are five shapes, five $F$ edges and six cherry
half turns.  The seven relations are the pentagon, three
cut-change naturality relations, and the two braid relations
and one disjoint-support relation.  Contracting an $F$ tree,
using the pentagon, and removing the three duplicate half
turns leaves exactly the three-generator Artin presentation
of $B_4$.  This proves (B.88) with its inverse maps.  Each
cherry half turn lies in a three-boundary sphere, each $F$
change lies in a four-boundary sphere, and each braid relation
is the three-strand isotopy in a four-boundary sphere.
The pentagon is the actual identity comparison of five cut
markings proved in B.14.  The seven relators thus come with
the specified geometric homotopies, not only word equalities.

Here are the covering functors themselves.  Write $U_m$ for
unordered configurations of $m$ points in the affine line.
The cubic family is
\[
 E_R:\ y^2=\prod_{r\in R}(x-r),\qquad |R|=3,
 \qquad p=\infty .                                             \tag{B.90}
\]
At infinity use the parameter $t=x/y$.  Translation of $x$
is contractible and does not change the leading tangent
coordinate.  After normalizing the cubic to have zero
quadratic coefficient, (B.90) gives the moduli of an elliptic
curve with its nonzero tangent vector at $p$.  Conversely,
the given tangent vector fixes the leading coefficients of
its Weierstrass functions $x,y$; translating $x$ to remove
the quadratic coefficient recovers its unordered branch
configuration.  The two constructions recover one another
up to the displayed coordinate changes.  Passing from
nonzero vectors to rays only removes a positive real,
contractible factor.  This is the cubic inverse construction
used in B.9, now with its actual framing retained.

Lift an adjacent positive half twist of two branch points
through the double cover.  In a disk containing exactly
those points the cover is an annulus.  The lift makes one
positive Dehn twist about its core; outside that disk it is
the identity.  Thus the two generators of the first row of
(B.88) go exactly to $a,b$.  The disk full twist
\[
 h=(ab)^3
\]
is the simultaneous once-around rotation of the three roots.
The scaling preserving the monic cubic has
$x\mapsto\lambda^2x$, $y\mapsto\lambda^3y$ and
$t\mapsto\lambda^{-1}t$.  A loop of $\lambda$ rotates the
roots twice and rotates the tangent frame once.  Hence the
positive elliptic full turn is $z=h^2$, exactly as in B.9.
This identifies the central lattice map, including its
factor two.

Let $\tau$ denote the positive half turn of this elliptic
boundary ray.  The full finite ray groupoid is the translation
groupoid on $\mathbb Z/2$ associated with
\[
 \left\langle B_3,\tau\ \middle|\
       \tau\text{ central},\ \tau^2=h^2\right\rangle
       \cong B_3\times C_2.                                   \tag{B.91}
\]
For the explicit isomorphism put $\omega=\tau h^{-1}$;
then $\omega^2=1$, and the inverse sends $\tau$ to $h\omega$.
The character to $\mathbb Z/2$ kills $a,b$ and sends $\omega$
to its nonzero element.  The deck involution of (B.90) sends
the tangent ray to its opposite and has square the identity.
It is natural in the branch configuration, so it commutes
with the braid paths.  Its composite with the positive
half rotation is the central loop whose square is $z=h^2$;
the positive orientation and the infinite cyclic centre
of $B_3$ identify that loop with $h$.  It therefore realizes
exactly $\omega=\tau h^{-1}$.  The positive boundary half turn is
$\tau=h\omega$, and its square is the required full turn.
Thus (B.91) supplies the actual half-turn functor and its
inverse, including their object maps.  Both sign objects
remain distinct and the sign action is free.

The cut change is equally explicit.  Put $s=aba$.  The
braid relation gives
\[
 sas^{-1}=b,\qquad sbs^{-1}=a,\qquad s^2=h.
                                                               \tag{B.92}
\]
It therefore sends the chosen nonseparating cut to its
once-intersecting mate.  Eliminating $b=a^{-1}sa^{-1}$ gives
the complete presentation
\[
 H_{1,1}=\langle a,s\mid (sa^{-1})^3=s^2\rangle,
       \qquad z=s^4.                                          \tag{B.93}
\]
The inverse substitutions are $a\mapsto a$ and $s\mapsto aba$.
Consequently all cut-change relations are checked in this
presentation; in the ray model one also has
$\tau=s^2\omega$, $\omega^2=1$, with $\omega$ central.
This records the genus-one cut change and its boundary
coherence on the same generators.

For two marked points the corresponding family is
\[
 E_R:\ y^2=\prod_{r\in R}(x-r),\qquad |R|=4,
       \qquad p_+,p_-\text{ over }\infty,                     \tag{B.94}
\]
where $y/x^2$ tends respectively to $+1,-1$ at the two
ordered points.  The inverse construction is the map from
$|p_++p_-|$: its quotient line, its affine coordinate up to
affine change, and the ordering of $p_+,p_-$ recover (B.94),
including the sign of $y$.  This is the inverse analytic-stack
construction in B.9.  Adjacent branch half twists lift to
the three-chain twists $a,b,c$.  At both points at infinity
$u=1/x$ is a local parameter.  The affine dilation loop
rotates the four roots once and has one tangent weight at
each point.  Its relation is therefore
\[
 H_{1,2}=(B_4\times\mathbb Z^2)
       /\langle (abc)^4=z_1z_2\rangle.                         \tag{B.95}
\]
This is the central induction of the disk full-twist lattice
along $1\mapsto(1,1)$.  The equivalence and inverse are the
maps just described in the family and on its coordinates.

The free ray model on $(\mathbb Z/2)^2$ is given without
suppressing either half turn:
\[
 \left\langle B_4,\tau_1,\tau_2\ \middle|\
 \tau_1,\tau_2\text{ central},\
       (abc)^4=\tau_1^2\tau_2^2\right\rangle,
 \quad\chi(\tau_i)=\overline e_i,
 \quad z_i=\tau_i^2.                                         \tag{B.96}
\]
The words of $B_4$ have character zero.  Therefore (B.96)
acts on its four sign objects as the actual positive half
rotations, and its isotropy is exactly (B.95).  The cover
functor carries these rotations to rotations of the two
coordinates $u$, so it respects composition and both full
turns.  No sign object is identified by the construction.

For completeness, all the local cut maps used in B.27 are
covered by this dictionary.  In the cubic case they are
$a$ and $h=s^2$ from (B.12).  In the quartic case use
\[
 \Delta=abcaba,\quad q=(abc)^2,\quad
 r=s^{-1}\Delta s,\quad h=(ab)^3.
\]
The identities
\[
 r^2=q^2=(abc)^4=z_1z_2,\qquad
 qaq^{-1}=c,\qquad h=ac^{-1}r
                                                               \tag{B.97}
\]
are precisely the branch exchanges and flag injections
of B.10--B.11.  The program checks these, (B.89)--(B.93),
and all full-twist centralities by the faithful Artin action.
The two-boundary cut changes $aba$ and $bcb$ satisfy their
three-strand relations in their displayed $B_3$ subgroups;
their fourth powers are the separating twists $(ab)^6$ and
$(bc)^6$.  Every word in these cuts and every incidence map
has the explicit substitution into $a,b,c,z_1,z_2$; the
complete presentation (B.95) proves that no additional
relation has been left unspecified at the groupoid level.

Finally, these are equivariant equivalences usable inside
B.27.  On a finite marking choose the $F$-tree paths of
(B.85), retain its sign coordinate, and apply the displayed
covering map.  The inverse replaces each $a,b,c$ by its
chosen $F/B$ word, and each half turn by (B.91) or (B.96).
The tree paths are the inverse natural isomorphism.  All
relations of that natural isomorphism are the already
listed naturality, pentagon, braid and central relations.
For a product of cut components apply the functors factor
by factor and then the explicit central pushout of B.1.
The resulting comparison respects the flag maps (B.97).
Transporting the general incidence maps through these
functors and their inverse natural isomorphisms gives the
same pseudofunctor coherence; B.27's strictification makes
it an ordinary finite diagram again.
\end{proof}
